% !PRINTED 569
\weberpage{569}
%<<CONT
\weq{20}{
\begin{aligned}
(-1,\eta)^2 &= 17 \\
(i,\eta)^2 &= (-1,\eta)(-1+4i) \\
(\alpha,\eta)^2 &= -(i,\eta)(3+i\sqrt{8});
\end{aligned}
}

ausserdem hat man $(1,\eta) = -1$,

\weq{21}{
\begin{aligned}
(i,\eta)(-i,\eta) &= 17 \\
(\alpha,\eta)(\alpha^{-1},\eta) &= 17 \\
(\alpha^3,\eta)(\alpha^{-3},\eta) &= 17
\end{aligned}
}

und nach (9)

\weq{22}{
(\alpha,\eta)(\alpha^3,\eta) = (-1,\eta)(3+i\sqrt{8}),
}

wodurch $(1,\eta)$, $(-1,\eta)$, $(\pm i,\eta)$, $(\alpha,\eta)$, $(\alpha^{-1},\eta)$, $(\alpha^{-3},\eta)$ durch
Quadratwurzeln bestimmt sind.
% [source-original mathematical omission, verified 2026-08-29] Formula (22), together
% with the already determined (alpha,eta), also determines (alpha^3,eta), but the printed
% enumeration omits that quantity while retaining its reciprocal partner (alpha^-3,eta).

Bei der Bestimmung, die wir über $r$ und $\alpha$ getroffen haben,
ergiebt sich leicht aus der Betrachtung der Werthe der Cosinus
(16), dass $(-1,\eta)$ positiv, also gleich der positiven Quadratwurzel
$\sqrt{17}$ ist, und dass ferner $(i,\eta)$ und $(\alpha,\eta)$ positive reelle
Theile haben, wodurch auch diese Grössen völlig bestimmt sind.
Es ist also

\[
\begin{aligned}
(-1,\eta) &= \sqrt{17} \\
(i,\eta) &= \sqrt[4]{17}\left\{\sqrt{\frac{\sqrt{17}-1}{2}} + i\sqrt{\frac{\sqrt{17}+1}{2}}\right\},
\end{aligned}
\]

wenn alle Wurzeln positiv genommen werden, und nach der
letzten Gleichung (20) lässt sich durch eine, wenn auch etwas
lange Formel $(\alpha,\eta)$ durch Quadratwurzeln aus reellen Grössen
darstellen. Auf die geometrischen Constructionen, die sich hier
anknüpfen, gehen wir nicht ein\footnote{Vgl. hierüber v. \gesperrt{Staudt}, Construction des regulären Siebzehnecks. Crelle's Journ. Bd. 24, 1842.}.

Die benutzte Vorzeichenbestimmung erhält man einfach,
wenn man nach (16) setzt

\[
\frac{1}{2}(-1,\eta) = \cos\frac{2\pi}{17} - \cos\frac{\pi}{17} + \cos\frac{8\pi}{17} + \cos\frac{4\pi}{17}
- \cos\frac{6\pi}{17} + \cos\frac{3\pi}{17} + \cos\frac{7\pi}{17} + \cos\frac{5\pi}{17}.
\]

Nach einer bekannten trigonometrischen Formel ist aber

\[
\begin{aligned}
\cos\frac{2\pi}{17} - \cos\frac{\pi}{17} &= -2\sin\frac{\pi}{34}\sin\frac{3\pi}{34} \\
-\cos\frac{6\pi}{17} + \cos\frac{5\pi}{17} &= 2\sin\frac{\pi}{34}\sin\frac{11\pi}{34},
\end{aligned}
\]
%CONT>>
